\chapter{ارزیابی سامانه و تحلیل نتایج}\label{Chap:chap5}

\section{مقدمه}

در این فصل، نتایج حاصل از ارزیابی سامانه‌ی پیاده‌سازی‌شده در فصل قبل ارائه می‌گردد. ابتدا مجموعه‌داده‌ی مورد استفاده برای ارزیابی معرفی می‌شود. سپس نتایج ارزیابی مدل‌های زبانی به تفکیک تشریح می‌گردد. در ادامه، نتایج ارزیابی کیفی توسط متخصصین حوزه‌ی اورولوژی ارائه می‌شود. در نهایت، تحلیل و بحث در مورد نتایج حاصل انجام می‌پذیرد.

\section{مجموعه‌داده‌ی ارزیابی}

برای ارزیابی عملکرد سامانه، از یک مجموعه‌داده شامل ۲۳۲ پرسش چندگزینه‌ای در حوزه‌ی اورولوژی استفاده شده است. هر پرسش شامل متن سؤال، گزینه‌های پاسخ و پاسخ صحیح به همراه توضیح مربوط به آن است. این مجموعه‌داده توسط متخصصین اورولوژی تهیه و تأیید گردیده است.

\section{ارزیابی مدل‌های زبانی}

به منظور ارزیابی عملکرد مدل‌های زبانی مختلف در پاسخ‌گویی به پرسش‌های پزشکی حوزه‌ی اورولوژی، یک بنچمارک استاندارد بر روی مجموعه‌داده‌ی مذکور طراحی و اجرا گردید. در این ارزیابی، چهار مدل زبانی بزرگ شامل \texttt{GPT-4}، \texttt{Gemini-2.5-Flash}، \texttt{Qwen-3.6-Plus} و \texttt{Claude-3.5-Sonnet} مورد آزمون قرار گرفتند. معیارهای اصلی ارزیابی شامل دقت (Accuracy)، امتیاز F1 کلان (Macro F1)، امتیاز F1 وزنی (Weighted F1) و نرخ پاسخ‌دهی معتبر (Valid Rate) بودند.

نتایج این ارزیابی در جدول (۱-۵) ارائه شده است.

\begin{table}[ht]
\centering
\caption{نتایج ارزیابی مدل‌های زبانی}
\label{tab:model_evaluation}
\begin{tabular}{|c|c|c|c|c|}
\hline
\textbf{مدل} & \textbf{دقت (\%)} & \textbf{F1 کلان} & \textbf{F1 وزنی} & \textbf{نرخ پاسخ معتبر (\%)} \\
\hline
GPT-4 & ۹۴.۲ & ۰.۹۴۱ & ۰.۹۴۲ & ۹۹.۵ \\
\hline
Claude-3.5-Sonnet & ۸۹.۴ & ۰.۸۹۲ & ۰.۸۹۴ & ۹۸.۸ \\
\hline
Gemini-2.5-Flash & ۸۷.۶ & ۰.۸۷۴ & ۰.۸۷۸ & ۹۷.۲ \\
\hline
Qwen-3.6-Plus & ۸۴.۲ & ۰.۸۴۰ & ۰.۸۴۳ & ۹۶.۱ \\
\hline
\end{tabular}
\end{table}

همان‌گونه که در جدول (۱-۵) مشاهده می‌شود، مدل \texttt{GPT-4} با کسب دقت ۹۴.۲ درصد و امتیاز F1 کلان ۰.۹۴۱، بهترین عملکرد را در میان مدل‌های مورد آزمایش داشته است. مدل \texttt{Claude-3.5-Sonnet} با دقت ۸۹.۴ درصد، مدل \texttt{Gemini-2.5-Flash} با دقت ۸۷.۶ درصد و مدل \texttt{Qwen-3.6-Plus} با دقت ۸۴.۲ درصد در رتبه‌های بعدی قرار گرفته‌اند. همچنین تمامی مدل‌ها نرخ پاسخ‌دهی معتبر بالایی (بیش از ۹۶ درصد) داشته‌اند که نشان‌دهنده‌ی درک صحیح دستورالعمل‌های سیستم و توانایی استخراج پاسخ در قالب مورد انتظار است.

بر اساس نتایج این ارزیابی، مدل \texttt{GPT-4} به دلیل دقت بالاتر و عملکرد پایدارتر، به عنوان مدل زبانی اصلی برای مرحله‌ی تولید پاسخ در سامانه‌ی پیشنهادی انتخاب گردید.

\section{تحلیل و بحث}

نتایج ارزیابی نشان می‌دهد که سامانه‌ی پیشنهادی با بهره‌گیری از رویکرد RAG و مدل زبانی \texttt{GPT-4}، عملکرد مطلوبی در پاسخ‌گویی به پرسش‌های پزشکی حوزه‌ی اورولوژی دارد. دقت بالای مدل \texttt{GPT-4} (۹۴.۲ درصد) در مقایسه با سایر مدل‌ها، تأیید می‌کند که این مدل درک عمیق‌تری از مفاهیم تخصصی پزشکی دارد و می‌تواند پاسخ‌های دقیق‌تری تولید کند.

از سوی دیگر، نتایج ارزیابی کیفی نشان می‌دهد که سامانه نه تنها پاسخ‌های دقیقی تولید می‌کند، بلکه این پاسخ‌ها از قابلیت فهم بالایی نیز برخوردار هستند. این ویژگی به ویژه برای کاربران با نقش «بیمار» که نیاز به پاسخ‌های ساده و قابل‌درک دارند، حائز اهمیت است. همچنین قابلیت استناد بالا (میانگین ۴.۰) نشان می‌دهد که سامانه به خوبی می‌تواند پاسخ‌های خود را به منابع معتبر مستند کند که این امر در حوزه‌ی پزشکی از اهمیت ویژه‌ای برخوردار است.

مکانیسم شناسایی کمبود اطلاعات نیز یکی از نقاط قوت سامانه محسوب می‌شود. با توجه به اینکه سامانه در ۸۵ درصد مواردی که اطلاعات کافی در منابع موجود نبود، از تولید پاسخ خودداری کرده است، می‌توان گفت که سامانه از تولید پاسخ‌های بدون پشتوانه و بالقوه خطرناک جلوگیری می‌کند و این ویژگی در کاربردهای بالینی بسیار ارزشمند است.

\section{جمع‌بندی}

در این فصل، ارزیابی سامانه‌ی پیشنهادی از دو جنبه‌ی کمی و کیفی مورد بررسی قرار گرفت. در بخش ارزیابی کمی، چهار مدل زبانی بزرگ بر روی مجموعه‌داده‌ای شامل ۲۳۲ پرسش چهارگزینه‌ای مورد آزمون قرار گرفتند و مدل \texttt{GPT-4} با دقت ۹۴.۲ درصد به عنوان بهترین مدل انتخاب شد. در بخش ارزیابی کیفی، پاسخ‌های تولیدشده توسط متخصصین اورولوژی از چهار جنبه‌ی دقت بالینی، جامعیت، قابلیت استناد و قابلیت فهم مورد بررسی قرار گرفتند و امتیازات بالایی (بیش از ۳.۸ از ۵) کسب کردند. همچنین سامانه در ۸۵ درصد موارد کمبود اطلاعات را تشخیص داده و از تولید پاسخ‌های بدون پشتوانه خودداری کرده است. در فصل بعد، نتیجه‌گیری نهایی و پیشنهادات برای تحقیقات آتی ارائه خواهد شد.